\section{Experimental Setup}
\label{sec:experimental_setup}

本节介绍本文对 \textit{Improving Code Localization with Repository Memory}
中方法的复现实验设定。该工作面向软件缺陷修复中的代码定位任务
（code localization），即给定一个 issue 描述和对应仓库版本，模型需要定位出
后续补丁中实际被修改的源代码文件。本文以 LocAgent 作为基础 baseline，
并进一步复现 Repository Memory（RepoMem）机制，分析其在代码定位任务中的有效性与局限性。

\subsection{Task Definition}
\label{subsec:task_definition}

给定一个代码仓库 $\mathcal{R}$、一个 issue 描述 $q$ 以及对应的
base commit $c_b$，代码定位任务要求模型输出一个候选文件列表
$\hat{F} = \{f_1, f_2, \ldots, f_k\}$。真实答案为该 issue 的补丁中被修改的文件集合
$F^\ast$。若真实修改文件集合 $F^\ast$ 被模型输出的 top-$k$ 文件完全覆盖，
则认为该样本在 Acc@$k$ 上定位成功：

\begin{equation}
    \text{Acc@}k =
    \mathbb{I}(F^\ast \subseteq \hat{F}_{1:k}),
\end{equation}

其中 $\mathbb{I}(\cdot)$ 为指示函数。本文主要报告 Acc@1、Acc@3 和 Acc@5。

\subsection{Dataset}
\label{subsec:dataset}

本文计划使用 SWE-bench Verified 作为主要复现实验数据集。SWE-bench Verified
由真实 GitHub issue 及其对应的 pull request 补丁构成，能够较好地评估模型在真实软件工程场景下的代码定位能力。

考虑到完整数据集复现成本较高，本文采用分阶段复现实验设置：

\begin{itemize}
    \item \textbf{Small setting:} 从 SWE-bench Verified 中选择 30 个样本，用于验证代码流程、评估脚本和 memory 构建流程是否正确。
    \item \textbf{Medium setting:} 在流程稳定后扩展到 50--100 个样本，用于获得更可靠的对比结果。
    \item \textbf{Full setting:} 若计算资源和 API 成本允许，进一步扩展到 SWE-bench Verified 完整测试集。
\end{itemize}

在初始复现实验中，本文优先选择样本数量较多、社区维护活跃的仓库，例如
\texttt{django}、\texttt{scikit-learn} 或 \texttt{sympy}，以便构建较充分的历史 commit memory。

\subsection{Baselines}
\label{subsec:baselines}

本文计划比较以下方法：

\begin{itemize}
    \item \textbf{LocAgent:}
    原始图引导代码定位 agent。该方法基于代码依赖图进行实体搜索和跳转，
    通过工具调用逐步定位可能需要修改的文件。LocAgent 不显式利用仓库历史信息，
    因此每个 issue 的定位过程基本从当前源码和 issue 描述开始。

    \item \textbf{RepoMem-Episodic:}
    仅使用 episodic memory 的 RepoMem 变体。该方法利用历史 commit message、
    修改文件、commit diff summary 以及可解析的 issue link，为 agent 提供历史修复经验。

    \item \textbf{RepoMem-Semantic:}
    仅使用 semantic memory 的 RepoMem 变体。该方法对仓库中历史上最活跃的文件生成自然语言摘要，
    帮助 agent 快速理解高频修改文件的职责和模块关系。

    \item \textbf{RepoMem-Full:}
    同时使用 episodic memory 和 semantic memory 的完整 RepoMem 设置。

    \item \textbf{Our Method:}
    在 RepoMem 基础上加入进一步改进。本文计划根据复现过程中的失败案例设计改进方法，
    例如动态 memory gating、time-aware commit retrieval 或 co-change graph memory。
\end{itemize}

\subsection{Repository Memory Construction}
\label{subsec:memory_construction}

对于每个测试样本，本文以其 base commit $c_b$ 作为当前仓库状态，只使用
$c_b$ 之前的历史信息构建 repository memory，避免使用未来信息造成数据泄漏。

\paragraph{Episodic Memory.}
对每个样本，本文计划收集 base commit 之前最多 $N=7000$ 条历史 commits。
每条 commit memory 包含以下字段：

\begin{itemize}
    \item commit hash；
    \item commit message；
    \item commit timestamp；
    \item changed files；
    \item diff 或 diff summary；
    \item commit message 中解析到的 issue id，如 \texttt{Fixes \#123} 或 \texttt{Closes \#123}。
\end{itemize}

本文使用 BM25 对 commit memory 建立检索索引，并提供以下两个工具：

\begin{itemize}
    \item \texttt{SearchCommit(query\_list, top\_k)}：根据 issue 描述或 agent 生成的查询检索相关历史 commits。
    \item \texttt{ExamineCommit(sha\_list, display\_issue)}：查看指定 commits 的详细信息，包括修改文件和 diff summary。
\end{itemize}

\paragraph{Semantic Memory.}
Semantic memory 用于总结仓库中高活跃文件的功能。本文首先统计历史 commits 中各文件被修改的频率，
选择 top-$M$ 个活跃文件。在 small setting 中，出于成本控制，本文计划设置 $M=50$；
在 medium/full setting 中，若资源允许，将扩展到 $M=200$。

对于每个活跃文件，使用大语言模型生成文件级摘要，摘要内容包括：

\begin{itemize}
    \item 文件的主要职责；
    \item 关键类、函数或数据结构；
    \item 与其他模块的关系；
    \item 该文件可能涉及的 bug 类型。
\end{itemize}

随后对 semantic summary 建立 BM25 索引，并提供以下两个工具：

\begin{itemize}
    \item \texttt{SearchSummary(query, top\_k)}：检索与当前 issue 相关的文件摘要。
    \item \texttt{ViewSummary(file)}：查看指定文件的完整 semantic summary。
\end{itemize}

\subsection{Implementation Details}
\label{subsec:implementation_details}

本文复现实验基于 LocAgent 官方实现进行扩展。首先复现原始 LocAgent 的代码定位流程，
包括仓库下载、dependency graph 构建、agent 推理和 file-level 评估。随后在不改变 LocAgent
原始代码工具的基础上，额外接入 RepoMem 的四个 memory 工具。

所有方法使用相同的数据划分、相同的 issue 输入和相同的评价脚本。为保证公平性，
不同方法使用相同的大语言模型作为 agent 后端。若实验中因 API 可用性或成本限制无法使用论文中的完全相同模型，
本文将在报告中明确说明模型版本差异。

实验过程中记录以下信息：

\begin{itemize}
    \item 每个样本的 predicted files；
    \item ground-truth modified files；
    \item Acc@1、Acc@3、Acc@5；
    \item agent 工具调用轨迹；
    \item token 使用量；
    \item API 成本；
    \item 运行时间；
    \item 失败样本的错误类型。
\end{itemize}

\subsection{Evaluation Metrics}
\label{subsec:evaluation_metrics}

本文主要使用 file-level localization accuracy 作为评价指标，包括 Acc@1、Acc@3 和 Acc@5。
其中 Acc@1 衡量模型最优先预测文件的准确性，Acc@3 和 Acc@5 衡量模型能否在较短候选列表中覆盖真实修改文件。

此外，为分析方法的实际使用成本，本文还计划统计：

\begin{itemize}
    \item \textbf{Average tool calls:} 每个样本平均工具调用次数；
    \item \textbf{Average tokens:} 每个样本平均输入和输出 token 数；
    \item \textbf{Average cost:} 每个样本平均 API 成本；
    \item \textbf{Average runtime:} 每个样本平均运行时间。
\end{itemize}

这些指标有助于分析 repository memory 是否不仅提升了定位准确率，同时也改善了 agent 的搜索效率。

\section{Experimental Plan}
\label{sec:experimental_plan}

本文复现实验分为四个阶段进行。

\subsection{Stage 1: Reproducing LocAgent}
\label{subsec:stage1}

第一阶段目标是跑通 LocAgent 原始方法，建立可靠的 baseline。具体步骤如下：

\begin{enumerate}
    \item 克隆 LocAgent 官方代码库并配置运行环境；
    \item 下载 SWE-bench Verified 数据集；
    \item 根据每个样本的 base commit 构建对应仓库状态；
    \item 构建 dependency graph；
    \item 在 small setting 上运行 LocAgent；
    \item 使用统一评估脚本计算 Acc@1、Acc@3 和 Acc@5。
\end{enumerate}

该阶段的预期产出包括 LocAgent 的小规模实验结果、运行日志和初步失败样本列表。
这一阶段的主要作用是确认实验流程是否正确，并为后续 RepoMem 复现提供对照结果。

\subsection{Stage 2: Reproducing Repository Memory}
\label{subsec:stage2}

第二阶段目标是构建 RepoMem 所需的 episodic memory 和 semantic memory。
具体步骤如下：

\begin{enumerate}
    \item 对每个样本收集 base commit 之前的历史 commits；
    \item 从 commit message 中解析 issue link；
    \item 提取 commit 修改文件和 diff summary；
    \item 使用 BM25 建立 episodic memory 检索索引；
    \item 根据历史修改频率选择 top active files；
    \item 使用大语言模型生成文件级 semantic summaries；
    \item 使用 BM25 建立 semantic memory 检索索引。
\end{enumerate}

该阶段的关键检查点是：对于任意一个 issue，agent 能否通过
\texttt{SearchCommit} 和 \texttt{SearchSummary} 检索到合理的历史信息。

\subsection{Stage 3: Integrating RepoMem with LocAgent}
\label{subsec:stage3}

第三阶段目标是将 repository memory 工具接入 LocAgent。
在该阶段中，本文保留 LocAgent 原有的代码搜索和图遍历工具，同时新增四个 memory 工具：

\begin{itemize}
    \item \texttt{SearchCommit}
    \item \texttt{ExamineCommit}
    \item \texttt{SearchSummary}
    \item \texttt{ViewSummary}
\end{itemize}

接入后，本文分别运行三种 RepoMem 设置：
RepoMem-Episodic、RepoMem-Semantic 和 RepoMem-Full。
通过与原始 LocAgent 对比，可以分析两类 memory 对代码定位性能的独立贡献。

\subsection{Stage 4: Analysis and Improvement}
\label{subsec:stage4}

第四阶段目标是进行结果分析，并基于失败案例提出本文的改进方法。
本文计划从以下角度分析 baseline 的失败原因：

\begin{itemize}
    \item LocAgent 是否因为缺少历史上下文而定位失败；
    \item episodic memory 是否能够召回与当前 issue 相关的历史修复；
    \item semantic memory 是否能够帮助 agent 理解高活跃文件的职责；
    \item RepoMem 是否会因为检索噪声被误导；
    \item memory 工具调用是否显著增加 token 和成本；
    \item agent 是否找到了相关文件但排序错误。
\end{itemize}

在失败分析基础上，本文计划提出一种改进方法。初步设想为
\textbf{GatedRepoMem}：在 agent 使用 repository memory 之前，先判断当前 issue
是否真的需要历史信息。如果 issue 涉及 regression、历史行为、兼容性或跨版本变化，
则增强 memory 检索；如果 issue 是局部 API、语法或简单逻辑错误，则降低 memory 使用频率，
避免无关历史 commits 干扰定位过程。

GatedRepoMem 的核心假设是：repository memory 并非对所有样本都有帮助。
在简单或高度局部的问题中，过多历史信息可能引入噪声；而在涉及历史行为和跨文件依赖的问题中，
memory 能够提供当前源码之外的关键线索。因此，通过动态决定 memory 的使用强度，
有望在保持 RepoMem 优势的同时降低噪声和成本。

\subsection{Expected Results}
\label{subsec:expected_results}

本文预期 LocAgent 在 small setting 上能够取得较强的 file-level localization 结果，
但在涉及历史行为、隐式模块关系和跨文件修改的样本中容易失败。
RepoMem-Episodic 预计能够改善历史修复模式明显的样本，
RepoMem-Semantic 预计能够帮助 agent 更快理解高活跃文件的职责。
RepoMem-Full 预期整体表现优于单独使用一种 memory 的变体。

对于本文提出的 GatedRepoMem，预期其主要优势不是在所有样本上显著提高 Acc@k，
而是在以下方面带来改进：

\begin{itemize}
    \item 减少无关 memory 检索造成的错误定位；
    \item 降低不必要的 token 和工具调用成本；
    \item 在 memory 有用的样本上保持 RepoMem 的定位优势；
    \item 提高 agent 工具使用轨迹的可解释性。
\end{itemize}

最终实验结果将通过主结果表、消融实验表和 case study 进行呈现。